\scanpage{15}

% archive sheet number: 14

Réciproquement, supposons $\psi$ \uline{un isomorphisme}, \uline{montrons qu'alors} $\varphi$ \uline{est un isomorphisme}, du moins si on admet que $C$ admet une base finie sur $A$. Soit en effet $\psi'$ l'isomorphisme inverse de $\psi$

\[
\psi':\operatorname{Hom}_{\unsure{A|B}}(C,C)\longrightarrow\operatorname{Hom}(B,C)
\]

il s'identifie à un hom

\[
\deleted{B\otimes}\operatorname{Hom}(C,C)\otimes B\longrightarrow C
\]

i.e. à un hom

\[
\psi_1^{\deleted{\unsure{1|*}}\prime}:B\longrightarrow\operatorname{Hom}(\operatorname{Hom}(C,C),C)
\]

On veut en tirer un hom $B\to C\otimes C$ en utilisant un hom naturel

\[
\unsure{\gamma|f|\rho}:\operatorname{Hom}(\operatorname{Hom}(C,C),C)\longrightarrow C\otimes C
\]

Pour construire un tel, on note que le premier mb est isomorphe à

\[
(C'\otimes C)'\otimes C=C\otimes C'\otimes C\deleted{=C\otimes\operatorname{Hom}(C,C)}
\]

\deleted{[en général, on a \deleted{$\operatorname{Hom}($} $P\otimes\operatorname{Hom}(Q,R)\to\operatorname{Hom}(\operatorname{Hom}(P,Q),R)$ i.e. $P\otimes\operatorname{Hom}(Q,R)\otimes\operatorname{Hom}(P,Q)\to R$, donné par $(p,v,u)\to(v\circ u)(p)$]},

or l'élément unité de $C^{\unsure{*|\times}}$ définit un $A$-hom naturel

\[
\varepsilon:C'\longrightarrow A\unsure{!|,}
\]

d'où un hom $\gamma=\mathrm{id}_C\otimes\varepsilon\otimes\mathrm{id}_C$.
